% =============================================================
% abstract.tex — Résumé en anglais, équivalent strict de resume.tex
% Composé à la suite du résumé, sur la même page (norme ESPRIT §H) :
% pas de \chapter*, qui provoquerait un saut de page ; le titre reprend
% le corps et la graisse des titres de chapitre.
% =============================================================
\addcontentsline{toc}{chapter}{Abstract}
\selectlanguage{english}

\vspace{10pt}
\noindent\hspace*{0.2cm}{\normalfont\bfseries\fontsize{18}{22}\selectfont Abstract}
\nopagebreak\par
\vspace{10pt}

{\footnotesize

This thesis presents the design, implementation, and evaluation of a secure Google Cloud hosting
foundation for MENAL-SARL. The starting point is a pilot application on a single server, without
full reproducibility or a security detection pipeline. The approach combines technical audit,
risk analysis, twenty verifiable requirements, a Zero Trust architecture, infrastructure as code
and blocking DevSecOps gates. The resulting foundation separates service identities, deploys
without permanent keys and runs seven detection rules over filtered log collection; a
domain-specific model maps alerts to MITRE ATT\&CK techniques, outside the incident score. On
19 August 2026, thirteen malicious requests from an end-to-end scenario were all blocked with
HTTP~403, found in the logs and aggregated by rule R2 fifteen minutes later. On 24 August 2026,
an unauthorised outbound connection and a cross-tenant connection were both refused and logged,
with matching timestamps on either side; the tenant data-isolation control passed all six
checks. Database restoration was measured at 32~min~45~s, with no data loss. The main
limitations are the ungoverned network egress of two workloads detached from the subnet, one
partially filtered path traversal pattern, the database instance still shared across tenants,
and incomplete cost and semantic quality measurements. The foundation is reproducible and
operable in the staging environment; its main contribution is not a new component but a method:
every control is paired with the test that demonstrates its effectiveness, and the report
publishes the controls that hold as well as those that do not hold yet.

\par\vspace{4pt}
\noindent\textbf{Keywords:} Zero Trust, DevSecOps, Google Cloud, infrastructure as code, SIEM,
MITRE ATT\&CK, multi-tenant security.

}

\selectlanguage{french}
